\chapter{String Tension, Planck Units, and the Heuristic Origin of Twenty-Six Dimensions}

These notes follow Leonard Susskind's fifth lecture on string theory, with visual curation from LazyingArt LLC and Video2Book. The lecture has a very definite two-part motion. We begin with an omission from earlier meetings, namely units, and use the old hadronic-string picture to isolate the one physical scale in the problem: the string tension. Only after that long preparatory arc through Planck units, direct detectability, gravity, and black holes does the lecture reset and turn to the older bosonic-string question of why the theory seems to want twenty-six spacetime dimensions.

\section{From Hooke's Law to String Tension}

We begin exactly where the lecture begins: with a stretched string treated as a spring. A nonrelativistic string, or if one likes an idealized rubber band, stores potential energy when its endpoints are pulled apart. The photographed board writes this in the familiar Hooke-law form.

\begin{figure}[t]
\centering
\includegraphics[width=0.72\textwidth]{lecture_05_figure_02.png}
\caption{Hooke-law potential for a stretched string. The board writes the endpoint separation as \(X\); in the cleaned derivation below we relabel it by \(\ell\).}
\end{figure}

Thus we write
\begin{equation}
E=\frac{k\ell^{2}}{2}.
\end{equation}
Here \(k\) is the spring constant and \(\ell\) is the distance between the endpoints. The lecture is careful about the factor of \(1/2\): it is missing at first and then restored.

The first nontrivial step is that this energy is not being interpreted as the mass of the string. In the infinite-momentum description, which has been lurking in the background from earlier lectures, the relevant internal nonrelativistic energy is associated with \(m^{2}\), not \(m\). In that sense,
\begin{equation}
m^{2}\propto k\ell^{2},
\end{equation}
so that, up to inessential factors such as \(\sqrt{2}\) that Susskind explicitly suppresses,
\begin{equation}
m\propto \sqrt{k}\,\ell.
\end{equation}

This immediately suggests a more natural parameter:
\begin{definition}
The string tension is
\begin{equation}
T=\sqrt{k}.
\end{equation}
It is the energy per unit length of the stretched string, and therefore also the force with which the string pulls back when we separate its endpoints.
\end{definition}

In the string rest frame this means simply that the energy of a stretched string grows linearly with its length,
\begin{equation}
E_{\text{rest}}\sim T\,\ell.
\end{equation}
This is the same logic as surface tension, except that now we are dealing with a one-dimensional object.

The lecture then pauses on units. In units with \(c=\hbar=1\), energy has the dimensions of inverse length, so
\begin{equation}
[E]=L^{-1},
\qquad
[T]=L^{-2}.
\end{equation}
Thus the tension has dimensions of energy squared, and that already tells us something important: it sets the scale of the theory in precisely the quantity in which string excitations are naturally measured, namely mass squared.

Susskind also reminds us that the same structure appears in the ordinary harmonic oscillator. For a spring of unit mass,
\begin{equation}
\omega\propto \sqrt{k},
\end{equation}
and the energy spacing between oscillator levels is proportional to \(\omega\). The stronger the spring, the larger the frequency; the larger the frequency, the larger the gap between successive levels. The string is no different in spirit. The tension is the quantity that controls the jump in mass squared from one excitation level to the next.

\section{Planck Units and the Scale of Fundamental Strings}

Once tension has been identified as the one scale, the lecture broadens from hadronic strings to hypothetical fundamental strings. If hadronic strings are already strong enough, in a laboratory fantasy, to support something like a truck, then fundamental strings are vastly stiffer still. This does not yet tell us the exact scale, but it tells us the direction: stiffer strings oscillate faster and fluctuate less, so their characteristic size is smaller and their excitation energy larger.

At this point the lecture asks what the Planck scale really means. The point is not merely numerical. The deeper question is how to choose units with genuinely universal content. A meter, a kilogram, and a second are convenient laboratory choices, but they are not fundamental. The lecture therefore chooses three universal dimensional constants:
\begin{equation}
[c]=LT^{-1},
\qquad
[\hbar]=ML^{2}T^{-1},
\qquad
[G]=L^{3}M^{-1}T^{-2}.
\end{equation}

We then ask for a combination
\begin{equation}
G^{p}\hbar^{q}c^{r}
\end{equation}
with dimensions of length squared. Matching mass, time, and length powers gives
\begin{align}
-p+q&=0,\\
-2p-q-r&=0,\\
3p+2q+r&=2.
\end{align}
The first equation gives \(q=p\); the second then gives \(r=-3p\); the third finally gives \(2p=2\), so \(p=1\). Therefore
\begin{equation}
\ell_{P}^{2}=\frac{G\hbar}{c^{3}}.
\end{equation}
This is the Planck area. Taking the square root gives the Planck length,
\begin{equation}
\ell_{P}=\sqrt{\frac{G\hbar}{c^{3}}},
\end{equation}
and from it we define the Planck time by the light-crossing time of a Planck-sized region,
\begin{equation}
t_{P}=\frac{\ell_{P}}{c}.
\end{equation}
The corresponding mass scale is
\begin{equation}
m_{P}=\sqrt{\frac{\hbar c}{G}}.
\end{equation}

The lecture keeps the numbers at the order-of-magnitude level:
\begin{equation}
\ell_{P}\sim 10^{-35}\,\text{m},
\qquad
t_{P}\sim 10^{-43}\,\text{s},
\qquad
m_{P}\sim 10^{-8}\,\text{kg}\sim 10^{19}\,\text{GeV}.
\end{equation}
The striking feature is that the length and time scales are absurdly small, while the mass scale is surprisingly macroscopic: a tiny dust grain in ordinary units.

This is the point at which the lecture makes its physical guess. The zero-point fluctuation size of a fundamental string is expected to be of order \(\ell_{P}\), and the energy required to excite one mode is expected to lie somewhere near the Planck energy.

\subsection{Question \& Answer}

\paragraph{Why think the string scale is near the Planck scale, and why not much smaller?}

The lecture's answer is deliberately cautious. One should not confuse this with numerology. The claim is not that the Planck length has been proved to be the string length; the claim is that it is the natural benchmark, and that known physics makes it hard to push the string scale very far away from it.

First, there is evidence that ordinary quantum field theory remains valid up to extremely short distances, perhaps only a few orders of magnitude above the Planck length. The standard model couplings appear, when extrapolated, to approach one another at a unification scale somewhat below the Planck scale. This does not pin down the string length exactly, but it suggests that stringy behavior does not set in at much larger distances.

Second, it is difficult to push the string scale below the Planck length. If we try to localize physics to a distance of order \(\ell_{P}\), the uncertainty principle tells us that the momentum uncertainty is of order
\begin{equation}
\Delta p\sim \frac{\hbar}{\ell_{P}},
\end{equation}
and therefore the associated energy is of order the Planck energy,
\begin{equation}
E\sim c\,\Delta p\sim m_{P}c^{2}.
\end{equation}
If we insist on making the characteristic size much smaller still, the required energy grows correspondingly. But now gravity becomes important. In the lecture's language, a very small object with very large mass is no longer naturally a string; it is pushed toward black-hole behavior. This is why string theorists are comfortable with a string scale somewhat larger than the Planck length, but not with one much smaller.

\section{Why Direct Tests Are So Hard}

Having constructed the Planck units, the lecture spends a long stretch turning them into phenomenology. The point is not to impress us with huge numbers for their own sake, but to make concrete what it would mean to excite a fundamental string.

Suppose an electron really were a string. Then to detect its stringy character directly we would want to do with the electron what hadron physics effectively does with the proton: hit it hard enough to set it vibrating and read off the excited spectrum. In the lecture's preferred normalization, the natural excitation scale is of order the Planck mass, or in particle-physics units
\begin{equation}
m_{P}\sim 10^{19}\,\text{GeV}.
\end{equation}
That is why direct excitation of an electron or photon looks hopeless. It is not merely that the total energy is large. One must concentrate that energy into an extremely small region.

The lecture also rephrases the gravitational question in the language of couplings. In electromagnetism the force law is schematically
\begin{equation}
F_{\mathrm{em}}\sim \frac{e^{2}}{R^{2}},
\end{equation}
whereas for gravitation it is
\begin{equation}
F_{\mathrm{grav}}\sim \frac{GM^{2}}{R^{2}}.
\end{equation}
Thus the dimensionless gravitational coupling at scale \(M\) behaves like
\begin{equation}
\alpha_{\mathrm{grav}}(M)\sim GM^{2}.
\end{equation}
Gravity is weak at low energies because \(M\) is small compared to the Planck mass. It becomes strong when \(M\) approaches \(m_{P}\).

This same scale governs black-hole production. The lecture distinguishes two mechanisms. One is ordinary gravitational collapse, in which a very heavy object such as a star pulls itself inward. The other is kinetic: collide sufficiently energetic particles in sufficiently small volume, and gravity can become strong enough to form a black hole. The first requires stars; the second requires Planck-scale collisions.

\begin{figure}[t]
\centering
\includegraphics[width=0.62\textwidth]{lecture_05_figure_03.png}
\caption{Susskind's qualitative energy sketch for gravitational collapse. The board is only partially legible, so it should be read as a schematic picture of competing branches and a collapse threshold, not as an exact graph.}
\end{figure}

The red board sketch used in this part of the lecture is qualitative rather than quantitative, but it captures the flavor of the discussion: different contributions compete until a threshold is reached, and beyond that threshold a steep branch rises sharply.

\begin{figure}[t]
\centering
\begin{tikzpicture}[scale=1.0]
\draw[->] (0,0) -- (0,4.2);
\draw[->] (0,0) -- (6.2,0);
\draw[dashed] (4.1,0) -- (4.1,3.0);
\fill (4.1,2.5) circle (0.07);
\draw (0.2,3.4) -- (4.1,2.5);
\draw (0.2,2.0) -- (4.1,2.0);
\draw (2.1,1.0) -- (4.1,2.5);
\draw (4.1,2.5) .. controls (4.6,2.7) and (5.0,3.0) .. (5.7,4.1);
\draw (4.6,0.7) .. controls (4.7,0.5) and (4.9,0.8) .. (5.0,0.6)
      .. controls (5.1,0.4) and (5.3,0.7) .. (5.4,0.5)
      .. controls (5.5,0.3) and (5.7,0.6) .. (5.8,0.4);
\node at (3.1,-0.35) {\small qualitative only};
\end{tikzpicture}
\caption{A cautious TikZ reconstruction of the qualitative structure of the sketch: several branches, a marked threshold, a dashed guide, and a steep right-hand rise. No exact scales or formulas are being asserted here.}
\end{figure}

The accelerator estimate is one of the lecture's most effective pieces of dimensional common sense. If a linear accelerator gives roughly \(100\) GeV in about two miles, then to reach \(10^{19}\) GeV one would need a machine larger by a factor of about \(10^{17}\). That pushes the required size into the galactic range. And once one remembers luminosity, the situation worsens. It is not enough to accelerate one particle to the Planck scale. One must collide enough particles often enough that some events probe Planckian distances. The lecture's conclusion is blunt: a really direct test is extraordinarily remote.

The discussion of extra dimensions is folded into the same logic. If extra dimensions were large compared to the string scale, ordinary quantum field theory would fail at energies where particles begin to explore them. Since field theory seems to work very well down to extremely short distances, the extra dimensions, if they exist, must be very small. And if one tries to make them much smaller than the Planck length, the same black-hole argument returns: the corresponding mass scale becomes too large, and the putative stringy or geometric object is driven into gravitational collapse.

\section{Reset: The Easy Version of Twenty-Six Dimensions}

After the break the lecture explicitly resets. We are told that the real mathematical machinery of string theory is conformal invariance, but that before entering that subject Susskind wants to give the \emph{easy version} of why bosonic string theory seems to require twenty-six dimensions. Easy, however, does not mean satisfying. The lecture presents it as historically suggestive and mathematically correct, but not as the deep explanation.

The starting point is the open string. We want the first excited state to be massless, because in the old bosonic story that state is interpreted as the photon. Therefore the ground state must sit one unit lower in mass squared:
\begin{equation}
m_{\text{ground}}^{2}=-1.
\end{equation}
Then one unit of excitation lifts the state to zero mass squared.

At this point the lecture explicitly sets aside a different problem, namely the fate of the negative-mass-squared ground state itself. For the moment we agree to keep it and ask only where the \(-1\) comes from.

The natural guess is that it must come from zero-point energy. Every oscillator contributes \(\frac{1}{2}\hbar\omega\), and in the lecture's units \(\hbar=1\). The photographed board captures the vacuum condition and its coordinate-space form.

\begin{figure}[t]
\centering
\includegraphics[width=0.70\textwidth]{lecture_05_figure_04.png}
\caption{Vacuum annihilation condition and ground-state equation. The board is partially obscured, but the operator condition and the Gaussian character of the solution are clearly visible.}
\end{figure}

In cleaned form the argument is the ordinary harmonic-oscillator one:
\begin{equation}
a^-|0\rangle=0.
\end{equation}
In coordinate space this becomes
\begin{equation}
\left(x-\frac{\partial}{\partial x}\right)\psi_{0}(x)=0,
\end{equation}
whose normalizable solution is
\begin{equation}
\psi_{0}(x)\propto e^{-x^{2}/2}.
\end{equation}
Thus the ground state is not at zero energy. It already contains the familiar zero-point contribution.

For the string, the \(n\)th mode has frequency
\begin{equation}
\omega_{n}=n,
\end{equation}
so the zero-point energy of that single oscillator is
\begin{equation}
E_{0,n}=\frac{1}{2}\omega_{n}=\frac{n}{2}.
\end{equation}
Therefore the formal ground-state energy is
\begin{equation}
E_{0}^{\text{formal}}=\frac{1}{2}\sum_{n=1}^{\infty}n.
\end{equation}
And here the tension of the lecture becomes explicit: all of these terms are positive. How could their sum possibly generate the negative shift that the photon requires?

\section{Light-Cone Energy and the Intercept Puzzle}

\subsection{Question \& Answer}

\paragraph{How can a sum of positive zero-point energies produce the negative ground-state shift needed for the photon?}

The lecture does not answer this by brute force. Instead it first backs up and says: we have been asking for the wrong quantity. In the infinite-momentum description, the physically interesting internal energy is not the total energy \(E\), because most of \(E\) is the trivial longitudinal motion of the whole system.

This is why the board next shows the large-\(P\) energy formula.

\begin{figure}[t]
\centering
\includegraphics[width=0.72\textwidth]{lecture_05_figure_05.png}
\caption{Subtracting longitudinal momentum from the energy. The board first writes the boosted energy and then isolates the internal piece by forming \(E-P\).}
\end{figure}

For a very fast system moving along the \(z\)-axis,
\begin{equation}
E=P+\frac{m^{2}}{2P}.
\end{equation}
The first term is just the overall conserved momentum. The second term contains the internal physics. Subtracting the trivial piece gives
\begin{equation}
E-P=\frac{m^{2}}{2P}.
\end{equation}
Multiplying by \(2P\) then recovers the quantity that we really care about:
\begin{equation}
2P(E-P)=m^{2}.
\end{equation}

This is the place where the lecture ties the old infinite-momentum discussion back to the present calculation. Internal excitations slow down under boosts by the usual time-dilation factor, so their raw energy differences become small as \(P\) grows. But once we strip off the trivial \(P\) piece and multiply back by \(2P\), what remains is the invariant mass squared.

This shift in viewpoint changes the zero-point problem. We do not literally need the raw sum of oscillator energies to equal \(-1\). It is enough if, after the correct infinite-momentum reinterpretation, the ground-state contribution looks like
\begin{equation}
-1+\text{(an infinite additive constant)}.
\end{equation}
The lecture insists that the infinite constant is not the interesting part. One already has such constant pieces in the boosted energy, and in this heuristic treatment they may be absorbed into them. What matters is whether a finite remainder survives once the divergent constant is separated off.

\section{Regulating the Zero-Point Sum}

Now the lecture turns to the mathematics. The board algebra at this point is lively and partly corrected in real time, so it is better to preserve the strategy cleanly than to imitate every intermediate line. The regulator is a small positive \(\epsilon\), inserted so that the sum converges before the limit \(\epsilon\to 0^{+}\) is taken.

\begin{figure}[t]
\centering
\includegraphics[width=0.68\textwidth]{lecture_05_figure_06.png}
\caption{Regulated sum and differentiation with respect to \(\epsilon\). The screenshot shows the essential move: turn the weighted sum into a derivative and display the first differentiated terms explicitly.}
\end{figure}

We begin with
\begin{equation}
S(\epsilon)=\sum_{n=1}^{\infty} n e^{-n\epsilon}.
\end{equation}
The first trick is to remove the factor of \(n\) by differentiating a simpler series:
\begin{equation}
S(\epsilon)
=
-\frac{\partial}{\partial\epsilon}\sum_{n=1}^{\infty} e^{-n\epsilon}.
\end{equation}
The sum inside is geometric:
\begin{equation}
\sum_{n=1}^{\infty} e^{-n\epsilon}
=
e^{-\epsilon}+e^{-2\epsilon}+e^{-3\epsilon}+\cdots
=
\frac{e^{-\epsilon}}{1-e^{-\epsilon}}.
\end{equation}
Hence
\begin{equation}
S(\epsilon)
=
-\frac{\partial}{\partial\epsilon}
\left(
\frac{e^{-\epsilon}}{1-e^{-\epsilon}}
\right).
\end{equation}

The lecture then expands for small \(\epsilon\). The cleaned result is
\begin{equation}
\frac{e^{-\epsilon}}{1-e^{-\epsilon}}
=
\frac{1}{\epsilon}-\frac{1}{2}+\frac{\epsilon}{12}+O(\epsilon^{3}),
\end{equation}
and therefore
\begin{equation}
S(\epsilon)
=
\frac{1}{\epsilon^{2}}-\frac{1}{12}+O(\epsilon^{2}).
\end{equation}

This is the central payoff of the regulated sum. It splits the formal divergence into two parts:
\begin{equation}
\sum_{n=1}^{\infty} n e^{-n\epsilon}
=
\underbrace{\frac{1}{\epsilon^{2}}}_{\text{divergent constant}}
\;+\;
\underbrace{\left(-\frac{1}{12}\right)}_{\text{finite part}}
\;+\;
O(\epsilon^{2}).
\end{equation}
The lecture treats the \(1/\epsilon^{2}\) term as an infinite additive constant, to be absorbed into a piece that is already present in the infinite-momentum energy bookkeeping. The finite piece is the famous
\begin{equation}
1+2+3+4+\cdots \rightsquigarrow -\frac{1}{12},
\end{equation}
where the right-hand side is understood not as an ordinary convergent sum but as the finite remainder extracted by this regularization.

At this point we restore the overall factor of \(1/2\) that comes from \(\frac{1}{2}\omega_{n}\). Thus one transverse oscillator contributes
\begin{equation}
E_{0}^{\text{per transverse oscillator}}=-\frac{1}{24}.
\end{equation}
This is the quantity that feeds into the dimension count.

\section{Counting Dimensions for Open and Closed Strings}

Now the end of the lecture comes quickly. Once one transverse oscillator contributes \(-1/24\), \(D_{\perp}\) transverse oscillation directions contribute
\begin{equation}
E_{0}^{\text{open}}=-\frac{D_{\perp}}{24}.
\end{equation}
For the open string we want the ground-state shift to be \(-1\), so that one quantum of excitation brings the state to zero mass squared. Therefore
\begin{equation}
-\frac{D_{\perp}}{24}=-1,
\end{equation}
and hence
\begin{equation}
D_{\perp}=24.
\end{equation}
Adding the longitudinal spatial direction and time gives
\begin{equation}
D=24+2=26.
\end{equation}

This is the lecture's heuristic explanation of why bosonic string theory comes with twenty-six spacetime dimensions. Historically, that was one of the first signs that something remarkable was happening. At the same time, the lecture repeatedly warns us not to overstate the argument. It is a real argument, but not the deep one.

The closed string then fits the same pattern. One now has both left-moving and right-moving oscillator sectors, so the contribution doubles:
\begin{equation}
E_{0}^{\text{closed}}
=
2\left(-\frac{D_{\perp}}{24}\right).
\end{equation}
With \(D_{\perp}=24\), this becomes
\begin{equation}
E_{0}^{\text{closed}}=-2.
\end{equation}
That is exactly what is needed in the lecture's normalization, because a massless spin-two state requires two units of excitation, one from each sector. The closed-string story therefore comes out consistently once the open-string count has fixed \(D_{\perp}=24\).

The chapter should end, as the lecture does, with the warning that this is not yet the real explanation. The deeper reason lies in conformal invariance on the worldsheet. That subject is deferred to the next lecture.

\section{Summary}

The lecture begins with a humble question about units and ends at the old bosonic answer \(D=26\). The bridge between the two is the string tension. A stretched string behaves like a spring, so its one physical parameter can be identified as a tension \(T\) with dimensions of energy squared. For fundamental strings that scale lies near the Planck regime, which is why direct tests look fantastically difficult and why gravity becomes unavoidable when one tries to push to shorter distances. After the break the lecture turns this same scale back onto the bosonic ground state: the zero-point energy of infinitely many oscillators is regulated, its finite remainder is extracted, and from the resulting \(-1/24\) per transverse oscillator one arrives at \(24\) transverse directions, hence \(26\) spacetime dimensions. The argument is historical, heuristic, and incomplete, but it is one of the clearest windows into how the old dimension count first appeared.